% !TEX program = pdflatex
% Proposition 1 page — deux articles JURIX 2026 — pour le directeur de thèse
\documentclass[fontsize=9pt,a4paper]{scrartcl}
\usepackage[T1]{fontenc}
\usepackage[utf8]{inputenc}
\usepackage[french]{babel}
\usepackage{newtxtext,newtxmath}
\renewcommand{\sfdefault}{qhv}       % sans-serif = TeX Gyre Heros (sans le package, évite le clash newtx)
\usepackage{microtype}
\usepackage{xcolor}
\usepackage[most]{tcolorbox}
\usepackage{enumitem}
\usepackage{fontawesome5}
\usepackage{tikz}
\usepackage[a4paper,top=1.0cm,bottom=0.9cm,left=1.2cm,right=1.2cm]{geometry}

% ── Palette (identité Pactiva : navy + or) ───────────────────────────────────
\definecolor{navy}{HTML}{12284B}
\definecolor{navsoft}{HTML}{2C4A73}
\definecolor{gold}{HTML}{B8892B}
\definecolor{ink}{HTML}{1A1F26}
\definecolor{surface}{HTML}{F4F6F9}
\definecolor{line}{HTML}{D7DEE7}

\color{ink}
\setlength{\parindent}{0pt}
\setlist[itemize]{leftmargin=1.15em,itemsep=1.5pt,topsep=2pt,parsep=0pt}
\renewcommand{\labelitemi}{\textcolor{gold}{\footnotesize\faAngleRight}}
\pagestyle{empty}

% ── Titre de section : icône + libellé sans-serif navy + filet or ────────────
\newcommand{\heading}[2]{%
  \par\vspace{7pt}%
  {\sffamily\bfseries\color{navy}\small #1\hspace{0.5em}#2}\par\vspace{1pt}%
  {\color{gold}\rule{\linewidth}{1pt}}\par\vspace{4pt}}

\newcommand{\badge}[1]{\tikz[baseline=-0.6ex]\node[fill=gold,text=white,font=\sffamily\bfseries\footnotesize,
  rounded corners=2pt,inner sep=2.2pt]{#1};}

\begin{document}

% ══ BANDEAU TITRE ════════════════════════════════════════════════════════════
\begin{tcolorbox}[enhanced,colback=navy,colframe=navy,boxrule=0pt,arc=3pt,
  left=12pt,right=12pt,top=8pt,bottom=8pt,
  borderline west={3pt}{0pt}{gold}]
  {\sffamily\bfseries\LARGE\color{white} JURIX 2026 — proposition de deux articles}\\[3pt]
  {\color{white!85}\small Valoriser l'annotation thématique des \emph{Terms of Service} et sa mise en
  graphe pour la \textbf{détection d'anomalies}}\\[5pt]
  {\sffamily\footnotesize\color{gold}Ahmed El Azhar Jebbari}\;{\color{white!70}\footnotesize\textbullet\;
  dir. Jean-Charles Lamirel \textbullet\; LORIA, Université de Lorraine \textbullet\; août 2026}
\end{tcolorbox}

\vspace{3pt}
% ── Fil rouge ────────────────────────────────────────────────────────────────
\begin{tcolorbox}[enhanced,colback=surface,colframe=line,boxrule=0.5pt,arc=2pt,
  left=8pt,right=8pt,top=5pt,bottom=5pt]
  {\sffamily\bfseries\color{navy}\footnotesize\faRoute~~Fil conducteur.}~
  {\small Le but de la thèse est de \textbf{détecter les anomalies} des contrats ; CLAUDETTE en est la
  base de référence reconnue (la \emph{clause abusive}, au sens de la directive 93/13/CEE, = l'anomalie).
  Nos \textbf{annotations thématiques multi-annotateurs} (clauses, thèmes, certitude) sont le
  \textbf{socle} qui rend cette détection possible, structurée et explicable. Les deux articles se suivent
  naturellement : \textbf{A produit et publie la ressource}, \textbf{B l'exploite pour détecter les
  anomalies} — ensemble, ils avancent le cœur de la thèse \emph{tout en} produisant deux publications.}
\end{tcolorbox}

% ══ JURIX EN BREF ════════════════════════════════════════════════════════════
\heading{\faUniversity}{JURIX en bref}
{\small
\begin{minipage}[t]{0.485\linewidth}
\begin{itemize}
  \item \textbf{Quoi} : conférence \textbf{européenne de référence en IA \& Droit} (\emph{Foundation for
  Legal KB Systems}) ; actes chez \textbf{IOS Press}, série FAIA.
  \item \textbf{Quand / où} : 39\textsuperscript{e} édition, \textbf{Toulouse (IRIT), 8–10 déc. 2026}.
  \item \textbf{Public} : communauté IA \& Droit, exigeante sur l'\textbf{ancrage juridique} et la
  \textbf{reproductibilité}.
\end{itemize}
\end{minipage}\hfill
\begin{minipage}[t]{0.485\linewidth}
\begin{itemize}
  \item \textbf{Soumission} : abstract \textbf{$\sim$28 août} \textbullet\ papier \textbf{$\sim$5 sept.}
  (EasyChair, \emph{single-blind} — pas d'anonymisation).
  \item \textbf{Formats} : long $\le$10\,p. ($\approx$20\,\% acc.) \textbullet\ short $\le$5\,p.
  \textbullet\ poster $\le$2\,p.
  \item \textbf{Choix stratégique} : \emph{short}/poster accessibles $\Rightarrow$ publication très
  probable pour un travail solide.
\end{itemize}
\end{minipage}}
\vspace{3pt}

{\small\textcolor{navy}{\faCheckCircle}~\textbf{Pourquoi nous y sommes chez nous :} la détection de
clauses abusives dans les ToS (\textbf{CLAUDETTE} : Lippi, Lagioia, Ruggeri, Galassi\ldots) est une
\textbf{lignée native de JURIX}. Notre travail s'y inscrit directement, avec un angle neuf.}

% ══ LES DEUX ARTICLES ════════════════════════════════════════════════════════
\heading{\faLayerGroup}{Deux articles complémentaires}

\begin{tcbraster}[raster columns=2,raster column skip=6pt,raster equal height=rows]
% ---- CARTE A ----
\begin{tcolorbox}[enhanced,colback=white,colframe=navy,boxrule=0.7pt,arc=2pt,
  left=6pt,right=6pt,top=5pt,bottom=5pt,
  title={\badge{A}\;\sffamily\bfseries\small\color{navy} Ressource \;\textcolor{navsoft}{\footnotesize(long)}},
  coltitle=navy,colbacktitle=white,fonttitle=\bfseries,
  attach title to upper=\hspace{0pt}\\[3pt]]
  {\sffamily\bfseries\footnotesize\color{ink}\emph{CLAUDETTE-Themes} — enrichissement thématique
  multi-annotateur d'Unfair-ToS}\par\vspace{3pt}
  {\sffamily\scriptsize\color{navy}\textbf{\faLightbulb~ABSTRACT (candidat)}}\par\vspace{1pt}
  {\footnotesize\itshape Nous enrichissons CLAUDETTE/UNFAIR-ToS : des étiquettes d'abusivité au niveau de
  la phrase, nous passons à une annotation thématique de clauses, \emph{multi-label} ($\approx$20 thèmes
  juridiques, certitude, justification), produite indépendamment par trois annotateurs formés, en anglais
  et en français. Nous mesurons l'accord inter-annotateur multi-label ($\alpha$ de Krippendorff, distance
  MASI, par thème) et publions un gold qui \emph{préserve le désaccord} — version arbitrée, votes bruts et
  soft labels — avec des baselines Legal-BERT. La ressource et son protocole de résolution de conflits
  sont réutilisables.}\par\vspace{3pt}
  {\sffamily\footnotesize\color{navy}\textbf{Contribution}}
  \begin{itemize}\footnotesize
    \item Comble une \textbf{limite reconnue} de CLAUDETTE (grain phrase\,$\to$\,clause, multi-label).
    \item \textbf{Ressource réutilisable} — JURIX valorise explicitement les \emph{datasets}.
    \item \textbf{Protocole de gold} transparent (cascade + désaccord préservé) = apport méthodologique.
  \end{itemize}
  \vspace{2pt}{\footnotesize\textcolor{gold}{\faDatabase}~\textbf{Valorise le travail d'annotation.}}
\end{tcolorbox}
% ---- CARTE B ----
\begin{tcolorbox}[enhanced,colback=white,colframe=gold,boxrule=0.7pt,arc=2pt,
  left=6pt,right=6pt,top=5pt,bottom=5pt,
  title={\badge{B}\;\sffamily\bfseries\small\color{navy} Graphe \;\textcolor{navsoft}{\footnotesize(short\,$\to$\,long)}},
  coltitle=navy,colbacktitle=white,fonttitle=\bfseries,
  attach title to upper=\hspace{0pt}\\[3pt]]
  {\sffamily\bfseries\footnotesize\color{ink}\emph{Hypergraph Co-Occurrence Anomalies} — découvrir les
  clauses abusives par graphe}\par\vspace{3pt}
  {\sffamily\scriptsize\color{navy}\textbf{\faLightbulb~ABSTRACT (candidat)}}\par\vspace{1pt}
  {\footnotesize\itshape Nous représentons chaque document de conditions d'utilisation comme un
  hypergraphe clause–thèmes, construit à partir des annotations thématiques multi-label, et posons la
  découverte de clauses abusives comme une détection d'anomalies de graphe \emph{non supervisée} : une
  combinaison de thèmes juridiques rare constitue une anomalie candidate. Nous évaluons les hyperarcs
  anormaux contre les étiquettes d'abusivité de CLAUDETTE (precision@k, AUC-PR) — et non par la seule
  rareté —, avec des baselines peu profondes et profondes, et mesurons par ablation le coût d'une couche
  déontique prédite (bruitée). Une détection relationnelle et explicable, bâtie sur la
  ressource~A.}\par\vspace{3pt}
  {\sffamily\footnotesize\color{navy}\textbf{Contribution}}
  \begin{itemize}\footnotesize
    \item Détection d'anomalies \textbf{relationnelle et explicable} (pas phrase par phrase).
    \item \textbf{Niche inédite} : hypergraphe multi-label sur ToS, ancré sur CLAUDETTE.
    \item \textbf{Construit sur la ressource~A} $\Rightarrow$ chaîne cohérente qui sert la thèse.
  \end{itemize}
  \vspace{2pt}{\footnotesize\textcolor{gold}{\faProjectDiagram}~\textbf{Valorise la détection d'anomalies.}}
\end{tcolorbox}
\end{tcbraster}

% ── Rigueur & positionnement ─────────────────────────────────────────────────
\vspace{5pt}
\begin{tcolorbox}[enhanced,colback=surface,colframe=line,boxrule=0.5pt,arc=2pt,
  left=8pt,right=8pt,top=5pt,bottom=5pt]
  {\sffamily\bfseries\color{navy}\footnotesize\faBalanceScale~~Rigueur \& positionnement.}~
  {\small Ancrage \textbf{protection du consommateur} (dir. 93/13/CEE) et dialogue avec la lignée
  CLAUDETTE (Lippi, Lagioia, Ruggeri, Galassi). Accord multi-label mesuré par \textbf{$\alpha$-MASI +
  Gwet AC} (robuste aux classes rares) ; \textbf{baselines} et \textbf{analyse d'erreurs} ; \textbf{désaccord
  préservé} (répond à Braun 2023, \emph{AI \& Law}). Reproductibilité : corpus, protocole et prompts
  publiés.}
\end{tcolorbox}

% ══ FAISABILITÉ AOÛT ═════════════════════════════════════════════════════════
\heading{\faClock}{Faisabilité en août \& ressources en place}

{\small Les données arrivent en août ; l'essentiel est ensuite de la \textbf{mise en forme} et de la
\textbf{rédaction}.}\par\vspace{4pt}

% mini-frise
{\centering\small
\textbf{mi-août} \emph{annotations 3\,$\times$\,50 ToS}
~\textcolor{gold}{\faLongArrowAltRight}~ \textbf{fin août} \emph{$\alpha$-MASI + gold + baselines}
~\textcolor{gold}{\faLongArrowAltRight}~ \textbf{$\sim$28 août} \emph{abstracts}
~\textcolor{gold}{\faLongArrowAltRight}~ \textbf{$\sim$5 sept.} \emph{papiers}\par}
\vspace{4pt}

{\small
\begin{minipage}[t]{0.485\linewidth}
{\sffamily\footnotesize\color{navy}\textbf{\faServer~~Déjà en place}}
\begin{itemize}
  \item Corpus \textbf{50 ToS} annotés, \textbf{EN + FR} (traduits phrase à phrase, alignés 1:1).
  \item Plateforme \textbf{Pactiva} + \textbf{moteur de gold en cascade} (arbitrage tracé) codé et testé.
  \item \textbf{Legal-BERT} (open-weights, \emph{sans LLM}) + code de segmentation/embeddings de la thèse.
  \item Métriques d'accord ($\alpha$-MASI, Gwet AC) et \textbf{export deux couches} outillés.
\end{itemize}
\end{minipage}\hfill
\begin{minipage}[t]{0.485\linewidth}
{\sffamily\footnotesize\color{navy}\textbf{\faLifeRing~~Le filet (dé-risque)}}
\begin{itemize}
  \item \textbf{A est sûr} : données présentes en août, méthodes sur étagère, surtout de la rédaction.
  \item \textbf{B n'est jamais bloquant} : soumis en \textbf{short} (preuve de concept) ou glissé en 2027.
  \item Arbitrage cadré sur un \textbf{échantillon stratifié} si le volume l'exige.
  \item \textbf{Compute modeste} (fine-tuning Legal-BERT + petit graphe), sans dépendance externe.
\end{itemize}
\end{minipage}}

% ══ CE QU'ON PROPOSE ═════════════════════════════════════════════════════════
\vspace{5pt}
\begin{tcolorbox}[enhanced,colback=navy,colframe=navy,boxrule=0pt,arc=2pt,
  left=10pt,right=10pt,top=6pt,bottom=6pt,borderline west={3pt}{0pt}{gold}]
  {\color{white}\small\textbf{\faFlag~~Proposition.} Soumettre à JURIX 2026 \textbf{A (long)} +
  \textbf{B (short)}. \textbf{A} valorise les \emph{thèmes} (le travail d'annotation), \textbf{B} les
  \emph{graphes} (la détection d'anomalies) ; les deux \textbf{avancent le cœur de la thèse}. \;
  \textbf{\faBullseye~Prochaine étape :} établir l'\textbf{état réel de la campagne} (nombre de ToS
  multi-annotés / arbitrés) pour verrouiller le périmètre et le calendrier.}
\end{tcolorbox}

\end{document}
